\section{Related Work}
\label{sec:related}

\para{Output-distribution uncertainty.}
A first family of detectors reads uncertainty from the output distribution.
Token likelihood and predictive entropy need one forward pass, while semantic entropy \citep{kuhn2023semantic} and SelfCheckGPT \citep{manakul2023selfcheckgpt} compare several sampled responses.
These methods use logits or samples, not attention.
We use token-likelihood statistics from the same forward pass as our main baseline, because they are the cheapest alternative to an attention-based signal.

\para{Hidden-state probes.}
A second family trains probes on internal activations.
SAPLMA trains a classifier on hidden states to predict truthfulness \citep{azaria2023internal}, INSIDE scores the consistency of internal representations across samples \citep{chen2024inside}, and \citet{orgad2025llms} study where correctness information sits in the residual stream.
Unlike these methods, we restrict the detector to a scalar summary of each attention row, which tests a specific and interpretable claim rather than whatever a probe can extract.

\para{Attention-based detectors.}
The closest work uses attention maps directly, but with statistics other than entropy.
Lookback Lens classifies contextual hallucinations in retrieval-augmented settings from the per-head ratio of attention on context versus generated tokens \citep{chuang2024lookback}.
LLM-Check scores attention kernels \citep{sriramanan2024llmcheck}.
\citet{yuksekgonul2024attention} relate factual errors to the attention mass placed on constraint tokens.
Building on this line, we isolate the simplest version of the idea, that attention is more diffuse when the model hallucinates, and test it with explicit length and domain controls.
Our finding that per-head features do much better than per-layer averages fits the head-level design of Lookback Lens.

\para{Data and evaluation.}
We draw questions from \popqa \citep{mallen2023popqa}, \sciq \citep{welbl2017sciq}, and \triviaqa \citep{joshi2017triviaqa}.
Because these sources differ in question style and in how often models answer correctly, we treat the source of a question as a confound to measure rather than ignore.

We did not run a systematic literature search for this study; this section covers the work we are aware of and is not exhaustive.
